\section{当前 LLM 评估的三大范式}

\subsection{三大评估范式概览}

当前评估 LLM 主要有三种方式：

\begin{enumerate}
    \item \textbf{困惑度}（Perplexity）：直接看训练/验证损失
    \item \textbf{综合基准平均}（Average over Everything）：在多个自动评估基准上取平均
    \item \textbf{竞技场式评估}（Arena-like）：模型间两两对比，由人类或 LLM 评判
\end{enumerate}

\begin{figure}[H]
\centering
\includegraphics[width=0.95\textwidth]{slides-images/slide-047.jpg}
\caption{当前 LLM 评估的三大范式：预训练模型通常展示困惑度和综合基准平均；微调模型通常展示综合基准平均和竞技场式评估\protect\footnotemark}
\end{figure}
\footnotetext{Slides 第48页。}

\begin{knowledgebox}{为什么预训练模型和微调模型使用不同的评估方式}
预训练模型主要展示困惑度和基准平均，因为困惑度在预训练阶段是直接可用的优化目标。微调模型则更多展示基准平均和竞技场式评估，因为微调后的模型其对数似然在目标数据集上不再校准（calibrated），困惑度不再是可靠的指标。
\end{knowledgebox}

\subsection{困惑度：简单但有效}

\begin{figure}[H]
\centering
\includegraphics[width=0.95\textwidth]{slides-images/slide-054.jpg}
\caption{困惑度与下游任务性能的高度相关性：x 轴为困惑度，y 轴为多个下游任务的平均分数。低困惑度的模型在下游任务上也表现更好\protect\footnotemark}
\end{figure}
\footnotetext{Slides 第55页。}

一个令人惊讶的发现：预训练困惑度（本质上就是``预测下一个词的能力''）与几乎所有下游任务的性能都高度相关。很多开发者在日常开发中只看困惑度，并信任它能代表下游性能。

\begin{warningbox}{困惑度的不可比性}
困惑度在以下两种情况下\textbf{不可比}：
\begin{enumerate}
    \item \textbf{不同数据集}：在不同数据集上计算的困惑度无法直接比较
    \item \textbf{不同分词器}：不同模型使用不同的分词器（tokenizer），词表大小不同。熵的上界是 $\log|\text{vocab}|$，词表大小为 1 时永远预测正确；词表大小为 100,000 时任务困难得多。因此 LLaMA 3 和 Gemini 即使在同一数据集上的困惑度也不可比
\end{enumerate}
\end{warningbox}

\subsection{综合基准平均：HELM 与 MMLU}

\begin{figure}[H]
\centering
\includegraphics[width=0.95\textwidth]{slides-images/slide-049.jpg}
\caption{常见的综合评估基准：GSM8K（小学数学）、MMLU（多任务语言理解）、LegalBench（法律）、MedQA（医学执照考试）等。HELM 和 HuggingFace Open LLM Leaderboard 汇聚了这些基准\protect\footnotemark}
\end{figure}
\footnotetext{Slides 第50页。}

\textbf{MMLU}（Massive Multitask Language Understanding）是目前最被广泛引用的单一基准：

\begin{itemize}
    \item 涵盖 \textbf{57 个不同领域}的多选题：形式逻辑、概念物理、计量经济学、高中生物学等
    \item 四选一的多选题格式，看似简单但评估实现细节很重要
    \item Mark Zuckerberg 发布 LLaMA 3 时都会引用 MMLU 分数
\end{itemize}

\begin{figure}[H]
\centering
\includegraphics[width=0.95\textwidth]{slides-images/slide-050.jpg}
\caption{MMLU 示例题目：天文学（超新星类型）和高中生物学（种群遗传学）的多选题\protect\footnotemark}
\end{figure}
\footnotetext{Slides 第51页。}

\textbf{代码评估}也越来越重要，原因有三：
\begin{enumerate}
    \item 代码能力与推理能力高度相关——代码写得好的模型通常推理也好
    \item 很多使用者关心编程辅助功能
    \item 代码评估天然容易自动化——写好单元测试，运行看是否通过即可
\end{enumerate}

\subsection{本章小结}

当前 LLM 评估的三大范式各有适用场景：困惑度适用于预训练阶段的快速评估（但跨模型不可比），综合基准平均（如 MMLU、HELM）适用于广泛能力的评估，竞技场式评估适用于对话/指令跟随能力的评估。实践中通常需要组合使用多种评估方式。

%% ========== Part 7: 评估中的问题与挑战 ==========

